% Copyright (c) 2024 by Daniel R. Driver.
% !TEX encoding = UTF-8 Unicode
% !TEX TS-program = XeLaTeX

\documentclass[titlepage]{article}

% This document presumes a file structure and set of inputs that are
% available at: git@github.com:danieldriver/syllabi.git

\newcommand\policy{../policy}
\newcommand\incl{../includes}
\ProvidesFile{variables.tex}[2022/09/11 v2.3 -- Syllabus variables]

\usepackage{xspace} % make manual spaces (like \mycmd\ ) unnecessary
\usepackage{xifthen} % provides \isempty test

% variables for internal use
\newcommand\prof{}
\newcommand\pdegree{}
\newcommand\pphone{}
\newcommand\pemail{}
\newcommand\poffice{}
\newcommand\phours{}
%
\newcommand\ccode{}
\newcommand\ctitle{}
\newcommand\cseries{}
\newcommand\cversion{}
\newcommand\csemester{}
\newcommand\cmeetson{}
\newcommand\cmeetsat{}
\newcommand\cmeetsin{}
\newcommand\cwebsite{}
\newcommand\cdescrip{}
\newcommand\cprereqs{}
\newcommand\edobject{}

% in case of fully online, BTh, or undergraduate courses - https://tex.stackexchange.com/a/5896
\newif\ifonline
\newcommand\Int[2]{\ifonline#1\else#2\fi}
\newif\ifbth
\newcommand\bth[2]{\ifbth#1\else#2\fi}
\newif\ifUG
\newcommand\UG[2]{\ifUG#1\else#2\fi}

% commands for setting variables in the preamble
\newcommand\professor[2][PhD]{
  \renewcommand\pdegree{#1\xspace}
  \renewcommand\prof{#2\xspace}}
\newcommand\phone[1]{
  \renewcommand\pphone{\addfontfeatures{Numbers=Monospaced}#1\xspace}}
\newcommand\email[1]{
  \renewcommand\pemail{\href{mailto:#1}{#1}\xspace}}
\newcommand\officehours[2][AST Library/Teams]{
  \renewcommand\poffice{#1\xspace}
  \renewcommand\phours{\href{https://outlook.office.com/bookwithme/user/b97b147a2e5c43b4a0e200c2cb630b77@astheology.ns.ca?anonymous}{#2}\xspace}%
  }
%
\newcommand\coursecode[2][1.0]{
  \renewcommand\cversion{#1\xspace}
  \renewcommand\ccode{#2\xspace}}
\newcommand\coursetitle[2][]{
  \ifthenelse{\isempty{#1}}%
    {}% do nothing if #1 is empty, else:
    {\renewcommand\cseries{#1\\[1ex]}}
  \renewcommand\ctitle{#2\xspace}}
\newcommand\semester[1]{
  \renewcommand\csemester{#1\xspace}}
\newcommand\meets[3]{
  \newcommand\AM{\textsc{am}}
  \newcommand\PM{\textsc{pm}}
  \renewcommand\cmeetson{#1\xspace}
  \renewcommand\cmeetsat{\Int{From 9:00 \AM}{#2}\xspace}
  \renewcommand\cmeetsin{\Int{\href{https://smu.brightspace.com/d2l/login}{Brightspace}}{#3}\xspace}}
\newcommand\website[2][]{
  \renewcommand\cwebsite{
  \ifthenelse{\isempty{#1}}%
    {\href{https://#2}}% copy url if #1 is empty, else
    {\href{#1}}% use specified url and print #2
  {#2}\xspace}}
\newcommand\cdescription[2][BF 1001; RM 1000 or GTRS 6000]{
  \renewcommand\cprereqs{#1}
  \renewcommand\cdescrip{#2\par}}
\newcommand\objectives[1]{
  \renewcommand\edobject{#1\par}}


%\onlinetrue % \Int{true}{false}
\coursecode[3.0]{HB 3116}
\coursetitle[Hebrew Bible]{Ecclesiastes}

% Taught as BSTH 4403 at Tyndale in:
%   - Fall 2015
% Taught as HB 3116 at AST in:
%   - Winter 2019
%   - Summer 2024

\professor{Daniel R. Driver}
\phone{902-425-7051}
\email{daniel.driver@astheology.ns.ca}
\officehours[AST Library, Main Floor]{Mon--Thur, 2:00--3:00 \PM}

\semester{Summer Term 2024}
\meets{Daily in Session 3}% \meets{on}{at}{in}
      {9:15 \AM--12:15 \PM}
      {Classroom 3}
\website[https://danieldriver.com/courses/hb-3116/]{danieldriver.com}
\cdescription{% copy from the Academic Calendar; add [] for unusual prereqs
Ecclesiastes will be studied in the context of the Writings and other
wisdom literature, and in connection with the book’s reception history.
By the end of the course students should be able to: understand the
basic shape of the Writings; have a critical appreciation of the idea of
wisdom literature; characterize some ancient and modern commentaries on
Ecclesiastes; develop their own readings, supporting them with detailed
exegesis.
}% end of course description
%\objectives{% recall Bloom's taxonomy: http://www.celt.iastate.edu/teaching/RevisedBlooms1.html
%
%In this course Ecclesiastes will be studied in the context of other
%biblical literature, with some reference to extra-biblical sources, and
%in connection with the book's reception history. By the end of the
%course students should be able to:
%	name biblical books classified with the Writings and wisdom literature;
%	understand some different ways this literature has been classified;
%	give examples of extra-biblical parallels to it;
%	locate Ecclesiastes relative to other Solomonic literature;
%	recognize quotes from Ecclesiastes by chapter;
%	name major ancient and modern commentators on Ecclesiastes;
%	report on specific commentators verbally;
%	articulate major interpretive options in reading Ecclesiastes;
%	defend their own preferred readings in writing, with detailed textual support.
%
%}% end of learning objectives

\ProvidesFile{preamble.tex}[2014/09/02 v1.1 -- Syllabus preamble]

% basic typography
\usepackage{fontspec,realscripts}
\defaultfontfeatures{Ligatures=TeX}
\setmainfont[%
  Numbers={Proportional,OldStyle},
  SmallCapsFeatures={RawFeature=+c2sc;+smcp,LetterSpace=1}]{Scala OT}
\newfontfamily\lining[Numbers={Lining}]{Scala OT}
\setsansfont[%
  SmallCapsFeatures={RawFeature=+c2sc;+smcp,LetterSpace=1}]{Scala Sans Pro}
\setmonofont[Scale=MatchLowercase]{Menlo}
\usepackage{sectsty}
\allsectionsfont{\sffamily\scshape}
\frenchspacing
\setlength{\emergencystretch}{3em} % prevent overfull lines

% custom font size and leading
\renewcommand\tiny{\fontsize{6}{9}\selectfont}
\renewcommand\scriptsize{\fontsize{7}{10}\selectfont}
\renewcommand\footnotesize{\fontsize{8}{10}\selectfont}
\renewcommand\small{\fontsize{9}{11}\selectfont}
\renewcommand\normalsize{\fontsize{10}{12}\selectfont}% base size
\renewcommand\large{\fontsize{11}{13}\selectfont}
\renewcommand\Large{\fontsize{12}{14}\selectfont}
\renewcommand\LARGE{\fontsize{16}{18}\selectfont}% "course syllabus \\ semester" benefits from more lead
\renewcommand\huge{\fontsize{20}{22}\selectfont}
\renewcommand\Huge{\fontsize{24}{26}\selectfont}

% layout packages: page, logo, tables
\usepackage[scale={0.6,0.8},
            xetex]{geometry}
\usepackage{graphicx}
\usepackage{array}     % allow insertions of column styling with >{}
\usepackage{booktabs}  % elegant horizontal rules in tables

% custom macros for a session count in the schedule of readings
\newcounter{session}
\newcounter{columns}
\newcounter{spancolumns}
\newcounter{midcolumns}
\newcounter{courseunit}
\newcommand\setcolumncount[2][0]{ % optionally set count to other than 0,
  \setcounter{session}{#1}        % e.g. to -1, or to a standing count
  \setcounter{columns}{#2}
  \setcounter{spancolumns}{#2-1}
  \setcounter{midcolumns}{#2-2}}
\newcommand\sessioncount{\stepcounter{session}\arabic{session}}
\newcommand\sessionskip[1]{\multicolumn{1}{@{}r@{ }}{#1}}
\newcommand\unit[1]{\multicolumn{\thecolumns}{c}{%
  \scshape\stepcounter{courseunit}\roman{courseunit}. #1}}
\newcommand\noclass[1]{\sessionskip{---.}&\multicolumn{\themidcolumns}{@{}l}{\itshape No Class: #1}}
\newcommand\rarr{\Large\char"EA31\hspace*{0.5pt}} % pointer finger w/ reminder tie in Scala OT
\newcommand\reminder[2]{\sessionskip{\rarr}&\multicolumn{\thespancolumns}{@{}l}{#1 \hfill #2}}

% color to match Tyndale's branding
\usepackage[usenames]{xcolor}
% predefined: black, white, red, green, blue, cyan, magenta, yellow
\definecolor{TyndaleURLs}{HTML}{0062A0} % links on tyndale.ca
\definecolor{TyndaleBlue}{cmyk}{1,1,0,.32}
\definecolor{TyndaleGold}{cmyk}{0,.27,1,0}
\definecolor{TyndaleRed}{cmyk}{0,1,.99,.04}
\definecolor{TyndaleBlack}{cmyk}{0,0,0,1}
\definecolor{TyndaleGreen}{cmyk}{.45,0,1,.24}
\definecolor{TyndaleOrange}{cmyk}{0,.79,1,0}
\definecolor{TyndaleAqua}{cmyk}{.47,0,.24,0}
\definecolor{TyndaleYellow}{cmyk}{.03,.03,.35,0}

% metadata (assumes a host of definitions are made in the main file)
\usepackage[setpagesize=false,     % leave this to geometry
            hyperfootnotes=false,  % fragile and distracting
            xetex]{hyperref}
\hypersetup{breaklinks=true,       % allow link text to break across lines
            colorlinks=true,       % colorlinks resets pdfborder to 0 0 0
            urlcolor=TyndaleURLs,  % for external links
            linkcolor=TyndaleRed,  % for normal internal links
            citecolor=TyndaleGold, % for bibliographical citations in text
            pdfauthor={\prof},
            pdftitle={\ccode: \ctitle},
            pdfsubject={Tyndale UC, \csemester},
            pdfcreator={github.com/danieldriver/syllabus}}
\urlstyle{same}                    % don't use monospace font for urls

% custom footlines
\usepackage{fancyhdr}
\pagestyle{fancy} % turn it on
\fancyhf{}        % reset everything
\renewcommand{\headrulewidth}{0pt} % remove header line as well
\lfoot{\sffamily\scshape\footnotesize\ctitle, \csemester}
\rfoot{\sffamily\scshape\footnotesize\prof\quad\thepage}

% language support: polyglossia needs to go last
\usepackage{polyglossia}
\setdefaultlanguage{english}
\setotherlanguage{hebrew}
\newfontfamily\hebrewfont[Script=Hebrew]{SBL BibLit}

% protect positioning of margin table in policy/grades -- needs to come after polyglossia
\usepackage{marginfix}

% gratuitous with custom title page, but useful as a fallback
\title{\ccode: \ctitle}
\author{\professor}
\date{\semester}


\begin{document}
\ProvidesFile{title.tex}[2013/09/06 v1.0 -- Syllabus title page]

\begin{titlepage}
  \begin{center}

    \begin{minipage}{\textwidth}
      \parbox[t]{0.5\textwidth}{
        \mbox{}\\[-\baselineskip] % dummy line to align parboxes
        \includegraphics[width=0.5\textwidth]{.syllabus/includes/TyndaleUC}}
      \hfill
      \parbox[t]{0.4\textwidth}{
        \raggedleft\LARGE\sffamily Course Syllabus\\
        \semester}
    \end{minipage}

    \vfill

    {\sffamily\textsc{\LARGE\MakeLowercase\coursecode}\\[2ex]
      \bfseries\itshape\Huge\coursetitle}

    \vfill

    \begin{tabular}{>{\bfseries}rl>{\bfseries}rl}
      \toprule
      Instructor & \professor                   & Course  & \liningnums{\coursecode}~\texttt{(v\version)} \\
      \midrule
      Phone      & \texttt{\phone}              & Meets   & \meeton                                       \\
      Email      & \href{mailto:\email}{\email} & Time    & \meetat                                       \\
      Office     & \officenumber                & Room    & \meetin                                       \\
      Hours      & \officehours                 & Website & \href{http://\website}{\website}              \\
      \bottomrule
    \end{tabular}

    \vfill

    \begin{description}
      \item[Commuter Hotline]
        Class cancellations due to inclement weather or illness will
        be announced on the commuter hotline at \texttt{416.226.6620
        x2187}. Alternately, weather cancellation information is posted
        at \href{http://tyndale.ca/weather}{tyndale.ca/weather}.
      \item[MyTyndale.ca]
        This course may have materials stored on its website, such as
        handouts or readings that may be needed in order to complete
        assignments. Students are responsible for checking these course
        pages on a regular basis. Here, too, students are able to view
        their grades throughout the semester. For more information see
        Section~\ref{mytyndale}, below.
      \item[Mail]
        Students are responsible for information communicated through
        their campus mailboxes and student e-mail accounts. A mailbox
        directory hangs beside the mailboxes. For more information
        contact the Registrar's office.
    \end{description}

  \end{center}

  \section{Course Description}
  \label{description}

  \emph{From the Academic Calendar:} \cdescription

\end{titlepage}
\setcounter{page}{2} % count the title page as page 1


%  \section{Learning Objectives}
%  \label{objectives}
%
%  \edobject

\section{Required Texts \& Materials}
\label{texts}

The following texts are required. Students are encouraged to obtain their own hard copies.

\begingroup
\renewcommand{\section}[2]{}% temporarily remove the section heading
\begin{thebibliography}{BS 1747–75}% use the longest item in the bibliography

	\bibitem[Qohelet]{Qohelet} Debra Band and Menachem Fisch.
    \emph{Qohelet: Searching for a Life Worth Living}.
    Waco, TX: Baylor University Press, 2023. AST Library copy on reserve.
%	ISBN 978-1481318730.

	\bibitem[TBD]{TBD} Select an additional commentary on Ecclesiates.
	You may find one upstairs in the library stacks (try BS 1747–1475),
	in a commentary kept in the reference section, or indeed somewhere
	else. Nominate your selection at the start of the second class. Be
	prepared to change commentaries if necessary.

\end{thebibliography}
\endgroup

Students may also want access to a good Study Bible. I recommend the OUP
editions of either the NRSV (ed. Michael Coogan) or NJPS (ed. Adele
Berlin and Marc Zvi Brettler).

\section{Supplementary Texts}
\label{supplementary}

Supplementary readings will be recommended throughout the course.
Excerpts from this literature may be placed on reserve or made available
electronically. The titles below, selected from AST’s library holdings,
are notable in relation to the aims of this course.

\begingroup
\renewcommand{\section}[2]{}% temporarily remove the section heading
\begin{thebibliography}{Christianson}% use the longest item in the bibliography

%	\bibitem[BHQ]{BHQ} \emph{Megilloth}. Vol. 18 of \emph{Biblia Hebraica Quinta}.
%	Edited by Adrian Schenker et al.
%	Stuttgart: Deutsche Bibelgesellschaft, 2004.

	\bibitem[Barbour]{Barbour} Barbour [Grillo], Jennie.
	\emph{The Story of Israel in the Book of Qohelet: Ecclesiastes as Cultural Memory}.
	Oxford: Oxford University Press, 2012. AST Library reserve: BS 1475.52 B37 2012.
%	See also Grillo, Jennifer.
%	ISBN 978-0199657827.

%	\bibitem[Bolin]{Bolin} Bolin, Thomas M.
%	\emph{Ecclesiastes and the Riddle of Authorship}.
%	London: Routledge, 2017.

	\bibitem[Christianson]{Christianson} Christianson, Eric S.
    \emph{Ecclesiastes Through the Centuries}.
    Malden, MA: Blackwell, 2012. \href{https://ast.primo.exlibrisgroup.com/view/action/uresolver.do?operation=resolveService&package_service_id=930405710007188&institutionId=7188&customerId=7185&VE=true}{AST eBook online}.
%	ISBN 978-0470674918.

%	\bibitem[Dell 2013]{Dell13} Dell, Katharine J.
%	\emph{Interpreting Ecclesiastes: Readers Old and New}.
%	Winona Lake, IN: Eisenbrauns, 2013.
%	ISBN 978-1575062815.

	\bibitem[Dell 2020]{Dell20} Dell, Katharine J.
	\emph{The Solomonic Corpus of “Wisdom” and Its Influence}.
	Oxford: Oxford University Press, 2020. \href{https://ast.primo.exlibrisgroup.com/view/action/uresolver.do?operation=resolveService&package_service_id=930405680007188&institutionId=7188&customerId=7185&VE=true}{AST eBook online}.
%	ISBN: 978-0198861560.

	\bibitem[Dell 2023]{Dell23} Dell, Katharine J.
	\emph{The Lord by Wisdom Founded the Earth: Creation and Covenant in Old Testament Theology}.
	Waco, TX: Baylor University Press, 2023. \href{https://ast.primo.exlibrisgroup.com/view/action/uresolver.do?operation=resolveService&package_service_id=930405660007188&institutionId=7188&customerId=7185&VE=true}{AST eBook online}.
%	ISBN 978-1481317047.

	\bibitem[Fox]{Fox} Fox, Michael V.
	\emph{Ecclesiastes: The Traditional Hebrew Text with the New JPS Translation}. The JPS Bible Commentary.
	Philadelphia: Jewish Publication Society, 2004. AST Library reserve: BS 1473 F69 2004.
%	ISBN 978-0827607422.

	\bibitem[Krüger]{Krüger} Krüger, Thomas.
	\emph{Qoheleth: A Commentary}. Translated by O. C. Dean Jr. Hermeneia: A Critical and Historical Commentary on the Bible.
	Minneapolis: Fortress, 2004. AST Library reserve: BS 1475.53 K78 2004.

	\bibitem[Short]{Short} Short, Robert L.
	\emph{A Time to Be Born -- A Time to Die}.
	New York: Harper \& Row, 1973. AST Library reserve: BS 1475.2 S56 1973; \href{https://archive.org/embed/timetobebornatim0000shor}{Internet Archive}.

	\bibitem[Stone]{Stone} Stone, Timothy J.
	\emph{The Compilational History of the Megilloth: Canon, Contoured Intertextuality and Meaning in the Writings}. FAT 2.59.
	Tübingen: Mohr Siebeck, 2013. \href{https://ast.primo.exlibrisgroup.com/view/action/uresolver.do?operation=resolveService&package_service_id=930405630007188&institutionId=7188&customerId=7185&VE=true}{AST eBook online}.
%	ISBN: 978-3161523755.

\end{thebibliography}
\endgroup

Also, these two general works are worth consulting and even owning.
%\cite{rlgs} in particular contains sound advice on core skills like
%reading religious texts, writing essays and reviews, revising essays,
%making oral presentations, and learning languages.

\begingroup
\renewcommand{\section}[2]{}% temporarily remove the section heading
\begin{thebibliography}{RLGS}% use the longest item in the bibliography

	\bibitem[RLGS]{rlgs} Northey, Margot, Bradford A. Anderson, and Joel N. Lohr.
	\emph{Making Sense in Religious Studies: A Student's Guide to Research and Writing}.
	3rd ed. Toronto: Oxford University Press, 2019. AST Library: Reference BL 41 N67 2019.

	\bibitem[SBL2]{sbl2} Collins, Billie Jean, et al.
	\emph{The SBL Handbook of Style}.
	2nd ed. Atlanta: SBL Press, 2014.
	Designed to augment \href{https://proxy.openathens.net/login?qurl=https%3A%2F%2Fwww.chicagomanualofstyle.org%2Fbook%2Fed17%2Ffrontmatter%2Ftoc.html}{\emph{Chicago Style}}
	(the standard at AST), there is also a free
	\href{https://www.sbl-site.org/assets/pdfs/pubs/SBLHSsupp2015-02.pdf}{Student Supplement for SBL2}. AST Library: Reference PN 147 S26 2014.

\end{thebibliography}
\endgroup

\section{Course Outline}
\label{outline}

We will adhere to the schedule in \autoref{schedule} as closely as
possible, though the professor reserves the right to adjust it to suit
the needs of the class.

\setcolumncount{4}% set up \sessioncount, \unit{}, \noclass{}, and \reminder{memo}{date} macros
\begin{table}[htb]% add p to put the schedule on its own page
  \centering
  \begin{tabular}{>{\sessioncount.}r@{ }llr}% make sure the column config agrees with \setcolumncount
	\toprule
	\sessionskip{\textbf{\S}.}&\textbf{Primary Text}&\textbf{Questions \emph{\&} Themes}&\textbf{Date}\\
	\midrule
	& Ecclesiastes 1     & What is \emph{hevel}?              & 15 July \\
	& Ecclesiastes 1--12 & Five questions of your own         & 16 July \\
	& Ecclesiastes 2     & Who is King Qohelet?               & 17 July \\
	& Ecclesiastes 3--4  & What time do you call this?        & 18 July \\
	& Ecclesiastes 5--6  & Can one find a perfect good?       & 19 July \\
	\noclass{Midterm Break}                                   & 20--21 July \\
	& Ecclesiastes 7--8  & Is that what they say?             & 22 July \\
	& Ecclesiastes 9--10 & What are the limits of knowledge?  & 23 July \\
	& Ecclesiastes 11    & How sweet is the light!            & 24 July \\
	& Ecclesiastes 12    & How does wisdom shape the book?    & 25 July \\
	& Ecclesiastes 1--12 & The sum of the matter              & 26 July \\
	\reminder{End of Term: Summer Administration hours begin}{29 July} \\
	\bottomrule
  \end{tabular}
  \caption{Schedule of Readings}
  \label{schedule}
\end{table}

See the AST website for a list of other \href{http://www.astheology.ns.ca/students/academic-dates.html}{important dates}.

\section{Evaluation}
\label{evaluation}

The grade structure for \ccode has the following elements.

\begin{enumerate}

	\item \textbf{Notes and quotes} will be solicited from students
	early in each class. These are to be drawn from the primary reading
	in \cite{Qohelet} and another commentary of your choice
	(\cite{TBD}). What do you note about the assigned material? What
	images, words, phrases, verses, or motifs stand out to you from text
	and translation, and from the visual and verbal commentary of Band
	and Fisch? What does your chosen commentator add? Make notes as you
	read. Be prepared to share some highlights with the class. Finally,
	be sure to address any particular prompts that the professor may set
	as homework for the day. Regular participation is expected and
	required.

	\item A \textbf{short paper} will give students an opportunity to
	work out an answer to a key question related to the interpretation
	of Ecclesiastes. The paper should be 1,500 words long, plus or minus
	10\%. Students may select from one of the following questions:

\begin{enumerate}
	\item How does the persona of King Qohelet emerge in Eccl 1–3?
	\item How should the term \emph{hevel} be translated in the book of Ecclesiastes?
	\item In what sense is Ecclesiastes a book of wisdom within the HB/OT?
	\item In what sense is Ecclesiastes rightly understood as a book of joy?
\end{enumerate}

	Alternatively, students may propose a different question, subject to
	approval by the professor. Final drafts of these papers are due at
	the start of class on Friday, 19 July, where they will be read aloud
	to a small group of peers. Students \emph{must} bring complete
	drafts with them to this workshop. Failure to do so will result in a
	loss of 50\% on the assignment grade. Revised final submissions are
	to be handed to the professor in hard copy, stapled on top of the
	workshop draft, at the start of class on Monday, 22 July.

	\item Students are to produce some form of \textbf{visual exegesis}
	related to a verse, short passage, or (at most) one chapter of
	Ecclesiastes. Band provides a superb example; others may be seen in
	\cite{Christianson} and \cite{Short}. These creative responses to
	the text will be presented to the class on the last two days of the
	courses (25–26 July). Students will retain possession of their own
	artwork. For evaluation, they will submit a 1,000 word essay
	explaining their art’s relation to the text of Ecclesiastes. For a
	model, look again to Band. The explanatory essays are due at the
	start of class on Friday, 26 July.

\end{enumerate}

The breakdown for the semester's total work is shown in
\autoref{grade-dist}.

\begin{table}[htbp]
  \centering
  {\lining
  \begin{tabular}{lr}
    \toprule
    Notes \& Quotes     & 30\% \\
    Short Essay         & 30\% \\
    Visual Exegesis     & 40\% \\
    \bottomrule
  \end{tabular}}
  \caption{Distribution of Grades}
  \label{grade-dist}
\end{table}

\ProvidesFile{grades.tex}[2023/09/06 v2.5 -- Course policy]

AST's \href{https://www.astheology.ns.ca/students/resources.html}{Academic
Calendar} provides guidelines and criteria for academic assessment. Marks are
assigned by letter grade using these benchmarks. Note that graduate and
undergraduate programs are currently assessed on a different scales at AST (see
\autoref{grade-syst}). For graduate students, a grade of C is deemed
``Unsatisfactory,'' whereas for undergraduate (BTh) students, a C is
``Satisfactory'' and the grade of D is possible for ``Marginal'' performance.
Other grades have different definitions, too.

\begin{table}[htbp]
 \centering
 {\lining\footnotesize
 \begin{tabular}{lll|ll}
   \toprule
               & \multicolumn{2}{c|}{Graduate Courses}   & \multicolumn{2}{c}{Undergraduate Courses} \\
               \cline{2-5}
   A+          & 94--100    & Exceptional               & 90--100    & Excellent                 \\
   A           & 87--93     & Outstanding               & 85--89     & Excellent                 \\
   A\char"2212 & 80--86     & Excellent                 & 80--84     & Excellent                 \\ [1ex]
   B+          & 77--79     & Good                      & 75--79     & Good                      \\
   B           & 73--76     & Acceptable                & 70--74     & Good                      \\
   B\char"2212 & 70--72     & Marginally Acceptable     &            &                           \\ [1ex]
   C           & 60--69     & Unsatisfactory            & 60--69     & Satisfactory              \\
   D           &            &                           & 50--59     & Marginal                  \\
   F           & 0--59      & Failure                   & 0--49      & Failure                   \\
   FP          & 0          & Failure due to Plagiarism & 0          & Failure due to Plagiarism \\
   \bottomrule
 \end{tabular}}
 \caption{Outline of Grades at AST}
 \label{grade-syst}
\end{table}

\ProvidesFile{other.tex}[2018/0/03 v2.4 -- Course policy]

\section{Other Course Policy}
\label{policy}

Late work will not be accepted, except in genuinely extenuating
circumstances. Students must submit something before the deadline if
they wish to receive credit. Unless I state otherwise, written
assignments are to be uploaded by 11:59 \PM\ on the date indicated.

Essay submissions must be typewritten, double-spaced, and formatted as
PDFs. They should be free from error. In this course they should follow
either \emph{The SBL Handbook of Style} (refer to the online
\href{https://www.sbl-site.org/assets/pdfs/pubs/SBLHSsupp2015-02.pdf}{
Student Supplement}), or \emph{The Chicago Manual of Style}, per AST's
policy for major papers. As a reminder, AST also upholds an Inclusive
Language Policy.

\href{http://www.eerdmans.com/Pages/Item/59043/Commentary-Statement.aspx}{Plagiarism},
if \href{https://www.theguardian.com/world/2013/feb/09/german-education-minister-quits-phd-plagiarism}{detected},
will result in failure of the course.

Students should request permission to record a class or lecture. If
permission is granted, or if recordings are provided (as in the case of
an online course), I stipulate that all recordings be for personal use
only. They may not be shared or distributed.

If you have abilities or disabilities that require modifications to the
assessment process or other aspects of this course, please advise the
course instructor as soon as possible.

Finally, I encourage the conscientious use of laptops, tablets, and
other technology in my classes. In classroom settings, realize that, as
\href{http://dx.doi.org/10.1016/j.compedu.2012.10.003}{cognitive
psychologists have demonstrated}, ``laptop multitasking hinders
classroom learning for both users and nearby peers.'' Do your part to
foster an environment of open dialogue by honouring the presence of your
classmates. In online settings, consider both the physical environment
in which you choose to work and the virtual environment that you help
create through your participation in various forums. Let your engagement
in this course be marked by rigour and charity alike.


\section{Bibliography}
\label{bib}

Articles and book chapters from the vast literature on Ecclesiastes in
reception, and in the context of the Writings and wisdom literature
(so-called), may be recommended as the semester progresses. Students are
encouraged to pursue secondary literature on their own as well. To start
your research, refer to the bibliographies in the books listed in
\autoref{supplementary}, above, search the library catalogue, and
consult relevant databases including:

\begin{itemize}
\item \href{https://www.oxfordbibliographies.com/display/document/obo-9780195393361/obo-9780195393361-0015.xml#obo-9780195393361-0015-div1-0005}{Oxford Bibliographies Online: Biblical Studies} -- s.v. “Ecclesiastes/Qohelet,” by Timothy J. Sandoval and Marcus Hayes;
\item \href{https://search.ebscohost.com/login.aspx?authtype=ip,shib&custid=s5315951&groupid=main&profile=ehost&defaultdb=ota}{Old Testament Abstracts};
\item \href{https://search.ebscohost.com/login.aspx?authtype=ip,shib&custid=s5315951&groupid=main&profile=ehost&defaultdb=rfh}{ATLA Religion Database (Full Text)};
\item \href{https://go.openathens.net/redirector/astheology.ns.ca?url=https://www.degruyter.com/document/doi/10.1515/ebr/html}{EBR: Encyclopedia of the Bible and Its Reception Online} --
s.v. “Ecclesiastes, Book of,” by various authors (and try other lemmas).
\end{itemize}

\end{document}
